% !TeX root = ../navier-stokes-manuscript.tex
% ============================================================
% 2.1 Vortex Stretching
% ============================================================

\subsection{A narrowing vortex against viscosity}\label{sub:ThePerfectlyStretchingVortex}

Follow one rotating structure inside the maximal velocity field \(u\).  Its
motion is still governed by the complete momentum equation
\[
  \partial_t u+(u\cdot\nabla)u=-\nabla p+\nu\Delta u.
\]
Transport carries the rotating fluid through its own spatial variation.  To see
whether that motion sharpens rather than merely relocates the structure, read
rotation and strain together.  Define
\[
  \omega:=\nabla\times u,
  \qquad
  S:=\frac12\bigl(\nabla u+(\nabla u)^{\mathsf T}\bigr).
\]
The vector \(\omega\) records local rotation, while \(S\) records how the surrounding velocity extends and contracts it.
When \(\omega\) aligns with an extensional direction of \(S\),
\[
  \omega\cdot S\omega>0,
\]
and the same strain amplifies the rotation.

That axial extension cannot occur alone.  Incompressibility requires
\[
  \operatorname{tr}S=\nabla\cdot u=0,
\]
so extension in one direction is paid for by contraction across the transverse section.
If the structure has width \(\ell\) and carries a velocity change \(\Delta U\) across it, then
\[
  |\nabla u|\sim\frac{\Delta U}{\ell}.
\]
The gradient grows only if \(\Delta U\) survives the loss of width.
This is the first place where the same object matters: a narrowing sketch is harmless unless the actual velocity field keeps a nontrivial change across it.

Kinetic energy does not prevent that possibility.
It measures
\[
  \int_{\R^3}|u(x,t)|^2\,dx,
\]
whereas the gradient compares the velocity change with the distance over which it occurs.
A structure can occupy less and less volume, carrying finite or even vanishing energy, while the ratio \(\Delta U/\ell\) diverges.

The equation must now decide whether the rotating structure can sustain this narrowing.
Taking the curl removes the pressure gradient because
\[
  \nabla\times\nabla p
  =0,
\]
and gives
\[
  \partial_t\omega+(u\cdot\nabla)\omega
  =
  (\omega\cdot\nabla)u+\nu\Delta\omega
  =S\omega+\nu\Delta\omega.
\]
The two terms on the right act simultaneously on the same vorticity: \(S\omega\) amplifies it through strain, while \(\nu\Delta\omega\) diffuses its spatial variation.
Their pointwise competition appears after pairing with \(\omega\):
\[
  \frac12(\partial_t+u\cdot\nabla)|\omega|^2
  =
  \omega\cdot S\omega
  +
  \nu\,\omega\cdot\Delta\omega.
\]
Aligned extension makes the first term positive.
The viscous term separates into a flux and a negative gradient contribution,
\[
  \nu\,\omega\cdot\Delta\omega
  =
  \frac{\nu}{2}\Delta|\omega|^2
  -
  \nu|\nabla\omega|^2.
\]
After integration over \(\R^3\), incompressibility removes transport and decay removes the Laplacian flux:
\[
  \frac12\frac{d}{dt}\norm{\omega(t)}_2^2
  +
  \nu\norm{\nabla\omega(t)}_2^2
  =
  \int_{\R^3}\omega\cdot S\omega\,dx.
\]
The balance is exact, but it has not chosen a winner.
When strain carries the rotation into a smaller width, stretching and diffusion both change with that width.
The next step is forced by this unresolved dependence: the entire coupled equation has to be compared across scale.

\divider{\NavierStokes{} Scaling}

Shrinking \(\ell\) by itself would change the object we are comparing.
The \NavierStokes{} equations tie width to velocity, pressure, and time, so a lawful scale change moves all four together.
For \(\lambda>0\), define
\[
  u_\lambda(x,t)
  :=
  \lambda u(\lambda x,\lambda^2t),
  \qquad
  p_\lambda(x,t)
  :=
  \lambda^2p(\lambda x,\lambda^2t).
\]
Substitution shows that every term acquires the same factor:
\[
\begin{aligned}
  \partial_t u_\lambda(x,t)
  &=
  \lambda^3(\partial_t u)(\lambda x,\lambda^2t),
  &
  (u_\lambda\cdot\nabla)u_\lambda(x,t)
  &=
  \lambda^3\bigl((u\cdot\nabla)u\bigr)(\lambda x,\lambda^2t),
  \\
  \nabla p_\lambda(x,t)
  &=
  \lambda^3(\nabla p)(\lambda x,\lambda^2t),
  &
  \Delta u_\lambda(x,t)
  &=
  \lambda^3(\Delta u)(\lambda x,\lambda^2t).
\end{aligned}
\]
Thus the rescaled field obeys the same balance, with
\[
  x\mapsto\lambda^{-1}x,
  \qquad
  u\mapsto\lambda u,
  \qquad
  p\mapsto\lambda^2p,
  \qquad
  t\mapsto\lambda^{-2}t.
\]
Moving to a smaller scale amplifies transport and viscosity by the same amount.
The equation therefore gives viscosity no automatic advantage merely because the structure is narrower.
This comparison relates rescaled solutions; it does not yet prove that one fixed fluid actually travels through those scales.

\divider{Critical Height}

We now need a quantity that reads the rescaled structure with the same size before and after this change.
The homogeneous Sobolev scale finds that height.
Let \(\Lambda:=(-\Delta)^{1/2}\).
For every \(s\in\R\) for which the homogeneous norm is finite,
\[
  \Lambda^s u_\lambda(x,t)
  =
  \lambda^{1+s}(\Lambda^s u)(\lambda x,\lambda^2t).
\]
Squaring and changing variables from \(x\) to \(y=\lambda x\) gives
\[
\begin{aligned}
  \norm{u_\lambda(t)}_{\dot H^s}^2
  &=
  \int_{\R^3}
  \lambda^{2+2s}
  \bigl|(\Lambda^s u)(\lambda x,\lambda^2t)\bigr|^2
  \,dx
  \\
  &=
  \lambda^{2s-1}
  \norm{u(\lambda^2t)}_{\dot H^s}^2.
\end{aligned}
\]
The scale factor disappears exactly when
\[
  2s-1=0,
  \qquad
  s=\frac12.
\]
Define the resulting critical height by
\[
  H(t)
  :=
  \frac12\norm{u(t)}_{\dot H^{1/2}}^2
  =
  \frac12\norm{\Lambda^{1/2}u(t)}_2^2.
\]

The energy law sits below this invariant level:
\begin{equation}\label{eq:BodyEnergyIdentity}
  \frac12\norm{u(t)}_2^2
  +
  \nu\int_0^t\norm{\nabla u(s)}_2^2\,ds
  =
  \frac12\norm{u_0}_2^2.
\end{equation}
Writing \(H_\lambda\) for the critical height of \(u_\lambda\), compare energy, critical height, and the first spatial derivative:
\[
  \norm{u_\lambda(t)}_2^2
  =
  \lambda^{-1}\norm{u(\lambda^2t)}_2^2,
  \qquad
  H_\lambda(t)=H(\lambda^2t),
  \qquad
  \norm{\nabla u_\lambda(t)}_2^2
  =
  \lambda\norm{\nabla u(\lambda^2t)}_2^2.
\]
As the structure is compressed, its energy reading falls, its critical height stays fixed, and its first-derivative reading rises.
This explains why finite energy can coexist with a sharpening flow: energy loses sight of the structure before the derivative does, while \(H\) continues to see it at full strength.

\divider{Nonlinear Production versus Viscous Dissipation}

Scale invariance has chosen the right gauge, but it has not said whether that gauge rises or falls along the actual history.
Apply \(\Lambda^{1/2}\) to the momentum equation and pair with \(\Lambda^{1/2}u\).
Keep transport and pressure together in the signed production
\[
  \mathcal P_{1/2}(t)
  :=
  -\operatorname{Re}
  \pair
    {\Lambda^{1/2}u(t)}
    {\Lambda^{1/2}\bigl((u(t)\cdot\nabla)u(t)+\nabla p(t)\bigr)}_{L^2}.
\]
The global pressure pairing vanishes after incompressibility is used, and viscosity contributes \(-\nu\norm{\Lambda^{3/2}u}_2^2\).
Hence
\[
  H'(t)
  +
  \nu\norm{\Lambda^{3/2}u(t)}_2^2
  =
  \mathcal P_{1/2}(t).
\]
The sign of \(\mathcal P_{1/2}\) is not prescribed.  Positive production raises
\(H\) only after it exceeds the viscous term beside it.
The original physical competition has now reached the one height that concentration cannot scale away.

\divider{The Heat Clock}

Scale does not hand viscosity the victory, but viscosity still assigns a clock to every width the flow reaches.
The linear heat equation makes that clock exact.
If \(v_t=\nu\Delta v\), then
\[
  \widehat v(\xi,t+\delta t)
  =
  e^{-\nu|\xi|^2\delta t}\widehat v(\xi,t).
\]
A structure of width \(r\) lives near frequency \(|\xi|\simeq r^{-1}\).
Its Fourier amplitude changes by order one once \(\nu|\xi|^2\delta t\) is of order one, so its viscous comparison time is
\begin{equation}\label{eq:BodyHeatComparisonTime}
  \tau_\nu(r)
  \asymp
  \frac{r^2}{\nu}.
\end{equation}
Halve the width and the time available to diffuse it falls by a factor of four.
Near a proposed terminal time \(T_*\), the width whose heat clock matches the time remaining is
\[
  r_\nu(t)
  \asymp
  \sqrt{\nu(T_*-t)}.
\]
This clock begins only after the nonlinear evolution has actually created and occupied the width.
The danger is therefore not a small scale by itself; it is a history able to reach the next small scale before the current clock runs out.

\divider{When the heat clocks fit before one time}

For geometrically shrinking widths, set
\[
  r_n:=2^{-n}r_0,
  \qquad
  \tau_n:=\frac{r_n^2}{\nu}.
\]
Using \eqref{eq:BodyHeatComparisonTime},
\[
  \tau_n
  =
  \frac{r_0^2}{\nu}4^{-n},
  \qquad
  \sum_{n=0}^{\infty}\tau_n
  =
  \frac{4r_0^2}{3\nu}
  <
  \infty.
\]
The clocks form a convergent series.  This is the Zeno arithmetic: infinitely
many scale changes can fit before one finite time.

For the Zeno picture to become a fluid singularity, one velocity history must form widths
\[
  r_0>r_1>r_2>\cdots\longrightarrow0
\]
at times
\[
  t_0<t_1<t_2<\cdots\uparrow T_*<\infty,
\]
with \(t_{n+1}-t_n\) following a summable schedule comparable to \(\tau_n\).
Each transition must occur while viscosity is already acting on the scale just formed.
The series proves only that the schedule fits; the coupled \NavierStokes{} evolution must still realize every step in the same fluid.

We do not need to assume that the fluid repeats one geometric profile.
If the maximal time were finite, the endpoint \(L^3\) criterion of Escauriaza, Seregin, and \v{S}ver\'ak would force \parencite{EscauriazaSereginSverak2003}
\[
  T_*<\infty
  \quad\Longrightarrow\quad
  \limsup_{t\uparrow T_*}\norm{u(t)}_{L^3}=\infty.
\]
But \(u=\Lambda^{-1/2}\Lambda^{1/2}u\), so the Hardy--Littlewood--Sobolev inequality gives \parencite[Theorem~4.3]{LiebLoss2001}
\[
  \norm{u(t)}_{L^3}
  \leq
  C\norm{u(t)}_{\dot H^{1/2}}
  =
  C\sqrt{2H(t)}.
\]
The critical height must therefore escape along some sequence \(t_n\uparrow T_*\):
\[
  H(t_n)\longrightarrow\infty.
\]
Plancherel and Cauchy--Schwarz then compare that escape with the first global derivative:
\[
\begin{aligned}
  2H(t)
  &=
  \norm{u(t)}_{\dot H^{1/2}}^2
  =
  \int_{\R^3}|\xi|\,|\widehat u(\xi,t)|^2\,d\xi
  \\
  &\leq
  \left(\int_{\R^3}|\widehat u(\xi,t)|^2\,d\xi\right)^{1/2}
  \left(\int_{\R^3}|\xi|^2|\widehat u(\xi,t)|^2\,d\xi\right)^{1/2}
  \\
  &=
  \norm{u(t)}_2\norm{\nabla u(t)}_2.
\end{aligned}
\]
The zero datum produces the zero solution and cannot lie on a finite-endpoint branch, so \(\norm{u_0}_2>0\) there.
Using the energy ceiling from \Cref{prop:ClassicalKineticEnergyIdentity} in the
last inequality gives
\[
  \norm{\nabla u(t_n)}_2
  \geq
  \frac{2H(t_n)}{\norm{u_0}_2}
  \longrightarrow
  \infty.
\]
Thus any finite last time forces both critical height and the global \(L^2\) gradient to become unbounded along a terminal sequence.
It does not require the pointwise velocity amplitude \(\norm{u(t)}_{L^\infty}\) to diverge.

One fixed-profile model shows how different readings of concentration can separate.
Choose a nonzero divergence-free profile \(\phi\in\Ssigma\) and set
\[
  u_\ell(x)
  :=
  \ell^{-1}\phi\!\left(\frac{x}{\ell}\right).
\]
Its amplitude, gradient, clock, and energy obey
\[
\begin{aligned}
  U_\ell
  :=
  \norm{u_\ell}_{L^\infty}
  &\sim
  \ell^{-1}
  \longrightarrow
  \infty,
  &
  \norm{\nabla u_\ell}_{L^\infty}
  &\sim
  \ell^{-2}
  \longrightarrow
  \infty,
  \\
  \tau_\ell
  &\asymp
  \frac{\ell^2}{\nu}
  \longrightarrow
  0,
  &
  E_\ell
  :=
  \frac12\norm{u_\ell}_2^2
  &\sim
  U_\ell^2\ell^3
  \sim
  \ell
  \longrightarrow
  0.
\end{aligned}
\]
Here amplitude and gradient diverge while the energy carried by the profile vanishes, yet
\[
  \frac12\norm{u_\ell}_{\dot H^{1/2}}^2
  =
  \frac12\norm{\phi}_{\dot H^{1/2}}^2,
\]
so critical height remains fixed.
This spike is an image of what scaling permits, not the general endpoint route above, where \(H(t_n)\) itself must diverge.

Both pictures return us to the same dynamical burden: one velocity history must move through sharper states on shorter clocks as \(t\uparrow T_*\).
Once the remaining time collapses, the next question is unavoidable: what must happen to the successive time responses of that same field?
